
% BEGIN S2
 
\par\medskip\noindent\textbf{2 The objects of interest}\par
 
% BEGIN S2.p1
 
Our development concerns discrete labelled Markov processes (LMPs) also called reactive probabilistic transition systems. They consist of a set of states $S$ and transitions $\tau_{a}:S\to\mathcal{D}S$ labelled with actions $a$ of some set $\mathcal{A}$, where $\mathcal{D}S$ is the set of discrete sub-distributions on $S$:\par
  
\par  \quad  \[\mathcal{D}S=\{\varphi\colon S\to[0,1]\mid\sum_{s\in S}\varphi(s)\leq 1\}.\]\quad  \quad  
  
Note that the size of the state set need not be restricted (for the sums to be defined), as the sum is defined as $\sum_{x\in S}\varphi(x)=\sup_{{\textrm{finite }}X\subseteq S}\left(\sum_{x\in X}\varphi(x)\right)$. It is easy to prove that discrete probability distributions always have a countable support, i.e., they assign non-zero probability to at most countably many elements of $S$.\par
 
% END S2.p1
 
% BEGIN S2.p2
 
For $s,t\in S$ and $a\in\mathcal{A}$, the value $\tau_{a}(s)(t)$ (also written $\tau_{a}(s,t)$) encodes the probability of jumping from $s$ to $t$ when action $a$ is taken. Figure 1 shows such a system with actions send and ack. The first transition is a subprobability: the missing probability represents the probability of failure. If $X\subseteq S$, we write $\tau_{a}(s)(X)$ for $\sum_{t\in X}\tau_{a}(s)(t)$. Coalgebraically, LMPs, for a set of actions or labels $\mathcal{A}$, are arrows $\tau\colon S\to(\mathcal{D}S)^{\mathcal{A}}$. We write $\tau_{a}(s)$ for the subdistribution $\tau(s)(a)$. In case of just one label, LMPs are (discrete time) Markov chains, i.e., coalgebras of the functor $\mathcal{D}$. The reason for working with subdistributions is that they come equipped with a non-trivial pointwise order. Bisimilarity of LMPs is the classical Larsen-Skou [38] bisimilarity.\par
 
% END S2.p2
 
% BEGIN S2.p3
 
We denote by $\mathbf{PMet_{1}}$ the category of $1$-bounded pseudometric spaces. It has as objects pseudometric spaces, which are pairs $\langle X,d\rangle$ of a set $X$ and a function $d\colon X\times X\to[0,1]$ which satisfies $d(x,x)=0$, is symmetric: $d(x,y)=d(y,x)$ and satisfies the triangular inequality: $d(x,y)+d(y,z)\geq d(x,z)$ for all $x,y,z\in X$. As arrows, $\mathbf{PMet_{1}}$ has non-expansive maps $f\colon\langle X,d_{X}\rangle\to\langle Y,d_{Y}\rangle$, i.e., functions $f\colon X\to Y$ satisfying $d_{Y}(f(x_{1}),f(x_{2}))\leq d_{X}(x_{1},x_{2})$. A metric or distance is a pseudometric that additionally satisfies $d(x,y)=0$ iff $x=y$. Of particular interest for LMPs are pseudometrics whose kernel is bisimilarity, that is, $d(x,y)=0$ iff $x$ and $y$ are bisimilar. As we only deal with pseudometrics here, we often use the word distance or metric as a shorthand for pseudometric. Following [59, 56], we will use the opposite pointwise order111One reason for this choice is to obtain the distance as the greatest, rather than least, fixpoint – just like bisimilarity is the greatest fixpoint of a suitable functional. on distances defined by:\par
  
\par  \quad  \[d_{1}\sqsubseteq d_{2}\quad\Longleftrightarrow\quad\forall x,y\in X:d_{1}(x,y)\geq d_{2}(x,y).\]\quad  \quad  
  
However, whenever it is clearer, we may use the direct pointwise order too: $d_{1}\geq d_{2}\Leftrightarrow d_{1}\sqsubseteq d_{2}$.\par
 
% END S2.p3
 
% BEGIN S2.SS1
 
\par\medskip\noindent\textbf{2.1 The Kantorovich pseudometric}\par
 
% BEGIN S2.SS1.p1
 
The observation that bisimulation is too strong in the context of probabilistic systems has led to the idea of defining a pseudometric that would give zero distance to bisimilar states. Initially, this pseudometric was defined using a real-valued logic. A set $\cal F$ of functionals were defined from states of LMPs to $[0,1]$, with the following syntax, mimicking logical formulas [24]:\par
  
\par  \quad  \[f:=1\mid\inf(f_{1},f_{2})\mid 1-f\mid f\varominus q\mid\langle a\rangle\,f\mbox{ with }q\in[0,1]\cap\mathbb{Q}.\]\quad  \quad  
  
We omit the semantics, but the next example will give a taste of it. A distance emerged naturally as\par
  
\par  \quad  \[d_{K}(s,t)=\sup_{f\in\cal F}\big|f(s)-f(t)\big|.\]\quad  \quad  
  
% END S2.SS1.p1
 
% BEGIN Thmtheorem1
 
\par\medskip\noindent\textbf{Example 1.}\par
 
% BEGIN Thmtheorem1.p1
 
As an example, we can look back at the channel of Figure 1. The maximum difference over functionals between $c_{\varepsilon}$ and $c_{\gamma}$, with $\varepsilon,\gamma>0$ is over the functional $\langle\textit{send\,}\rangle\langle\textit{ack}\,\rangle 1$, which evaluates to $1-\varepsilon$ and $1-\gamma$, respectively, and thus yields the distance $d_{K}(c_{\varepsilon},c_{\gamma})=|\varepsilon-\gamma|$. Between $c_{\varepsilon}$ and $c_{0}$, the functional $(\langle send\rangle\langle ack\rangle)^{n}1$ evaluates to $(1-\varepsilon)^{n}$ on $c_{\varepsilon}$, and to 1 on $c_{0}$. So the supremum results in $d_{K}(c_{\varepsilon},c_{0})=1$.\par
 
% END Thmtheorem1.p1
 
% END Thmtheorem1
 
% BEGIN S2.SS1.p2
 
Later on, it was proven that the functionals could be any non-expansive maps and that this distance is the distance obtained via a final coalgebra construction for a functor (monad) that is a lifting of the Distribution functor (monad) to pseudometric spaces using the Kantorovich metric, by van Breugel and Worrell [59] and van Breugel et al. [56, 57].\par
 
% END S2.SS1.p2
 
% BEGIN S2.SS1.p3
 
This distance is a fixpoint of the functional on distances on the states of an LMP $\tau:S\to(\mathcal{D}S)^{\mathcal{A}}$\par
  
\par  \quad  \[\Delta_{K}(d)(s_{0},s_{1})=\sup_{a\in\mathcal{A}}\,\delta_{K}^{d}(\tau_{a}(s_{0}),\tau_{a}(s_{1}))\]\quad  \quad  (1)\quad  
  
where $\delta_{K}^{d}$ is the well-known Kantorovich distance between two distributions. We do not need its definition here, the interested reader can find it in, e.g., [56].\par
 
% END S2.SS1.p3
  
% BEGIN S2.SS1.p4
 
The distance on distributions is just a parameter in Equation (1); changing the Kantorovich distance $\delta_{K}^{d}$ to another distance on distributions and looking for a fixpoint will give us a new distance on states. In fact we are interested in the question formulated the other way around: find a distance on distributions for which the fixpoint is the $\varepsilon$-distance $d^{*}$, defined in the next section. Similarly, the functional $\Delta_{K}$ was discovered after the behavioural distance $d_{K}$ was introduced. However, it could have been done the other way around, choosing a distance on distributions, and looking for the resulting distance on states, as is done by Baldan et al. [9]. We discuss this in more detail again in Section 3 below.\par
 
% END S2.SS1.p4
 
% END S2.SS1
 
% BEGIN S2.SS2
 
\par\medskip\noindent\textbf{2.2 $\varepsilon$-Bisimulation}\par
 
% BEGIN S2.SS2.p1
 
Bisimulation, being too strong for the comparison of states in quantitative systems, has also led to relaxing the definition of simulation and bisimulation itself to approximate relations. As mentioned in the introduction, there are a few of those definitions of approximate bisimulations, but we are interested in the following.\par
 
% END S2.SS2.p1
 
% BEGIN Thmtheorem2
 
\par\medskip\noindent\textbf{Definition 2 ([26]).}\par
 
% BEGIN Thmtheorem2.p1
 
Let $\tau:S\to(\mathcal{D}S)^{\mathcal{A}}$ be an LMP and let $\varepsilon\in[0,1]$. A relation $R\subseteq S\times S$ is an $\varepsilon$-simulation if whenever $sRt$, then for all $a\in\mathcal{A}$, for all $X\subseteq S$\par
  
\par  \quad  \[\tau_{a}(s)(X)\leq\tau_{a}(t)({R}(X))+\varepsilon,\qquad\mbox{ where }R(X)=\{y\mid\exists x\in X:xRy\}.\]\quad  \quad  
  
If $R$ is symmetric, it is an $\varepsilon$-bisimulation. A state $s$ is $\varepsilon$-simulated by state $t$, written $s\prec_{\varepsilon}t$, if $sRt$ for some $\varepsilon$-simulation $R$. If $sRt$ for $R$ $\varepsilon$-bisimulation, we write $s\sim_{\varepsilon}t$ and we say that $s$ and $t$ are $\varepsilon$-bisimilar.\par
 
% END Thmtheorem2.p1
 
% END Thmtheorem2
 
% BEGIN S2.SS2.p2
 
As expected, ordinary (bi)simulation on LMPs [38, 11] is simply $0$-(bi)simulation. This definition has an extension to nondeterministic and probabilistic finite systems [54].\par
 
% END S2.SS2.p2
 
% BEGIN S2.SS2.p3
 
The operation $R(X)$ on a set $X$ is what restricts this work to discrete distributions. In general, if $X$ is measurable, we may not have $R(X)$ measurable. Working in analytic spaces, as was done before for LMPs [23] may solve this issue, but we leave it for future work. In particular, one could ask $R(X)$ to be measurable whenever $X$ is, as done in [10], but many results are not proven in that case, like the logical characterisation, or Theorem 5 below.\par
 
% END S2.SS2.p3
 
% BEGIN Thmtheorem3
 
\par\medskip\noindent\textbf{Example 3.}\par
 
% BEGIN Thmtheorem3.p1
 
Consider the following example, with $\gamma\in(0,1]$.\par
 
% END Thmtheorem3.p1
 
% BEGIN Thmtheorem3.p2
 
\begin{center}\includegraphics[width=0.95\linewidth,height=0.45\textheight,keepaspectratio]{graphic-02.pdf}\end{center}
 
% END Thmtheorem3.p2
 
% BEGIN Thmtheorem3.p3
 
The relation $R=\{(s,t),(s_{1},t_{1})\}$ is an $\varepsilon$-simulation for $\varepsilon=\gamma$, and so is $R\cup R^{-1}$. Hence, $s\prec_{{\gamma}}t$, $t\prec_{{\gamma}}s$ and $s\sim_{{\gamma}}t$. However, the situation is different for $\varepsilon<\gamma$: for example, taking $\varepsilon=0$, we observe that $s\prec_{0}t$ (with the relation $R\cup\{(s_{2},t_{1})\}$). However, $t\not\prec_{0}s$, because a relation $R^{\prime}$ relating $t$ and $s$ would have to also relate $t_{1}$ to $s_{1}$ and to $s_{2}$, if $\gamma<1$, and $t_{1}$ to $s_{2}$ if $\gamma=1$. Indeed, taking $X=\{t_{1}\}$, we need $\tau_{a}(t)(X)\leq\tau_{a}(s)(R^{\prime}(X))+0$, but for this to be the case, we would need $\tau_{a}(s)(R^{\prime}(X))=1$ and hence $R^{\prime}$ must include the pairs $(t_{1},s_{1})$ and $(t_{1},s_{2})$ if $\gamma<1$, and otherwise must include $(t_{1},s_{2})$ if $\gamma=1$. So, in any case $(t_{1},s_{2})\in R^{\prime}$, but the pair $(t_{1},s_{2})$ cannot be related by any $\varepsilon$-simulation relation $R^{\prime}$ for $\varepsilon<1$ (which is the case here as $\varepsilon=0$), as, for example, taking again $X=\{t_{1}\}$, we would have $\tau_{a}(t_{1})(X)=1$ but $\tau_{a}(s_{2})(R^{\prime}(X))=0$. Hence, also $s\not\sim_{0}t$.\par
 
% END Thmtheorem3.p3
 
% END Thmtheorem3
 
% BEGIN S2.SS2.p4
 
In contrast to the case of 0-bisimilarity, two-way $\varepsilon$-similarity is not $\varepsilon$-bisimilarity [26, 54]. From the notion of $\varepsilon$-bisimilarity, a pseudometric arises naturally by taking the infimum.\par
 
% END S2.SS2.p4
 
% BEGIN Thmtheorem4
 
\par\medskip\noindent\textbf{Definition 4 (The $\varepsilon$-distance $d^{*}$ [26]).}\par
 
% BEGIN Thmtheorem4.p1
 
Let $\tau:S\to(\mathcal{D}S)^{\mathcal{A}}$ be an LMP. The pseudometric $d^{*}$ on $S\times S$, called $\varepsilon$-distance, is defined as follows:\par
  
\par  \quad  \[\begin{array}[]{lcl}d^{*}:&S\times S&\rightarrow[0,1]\\
&(s,t)&\mapsto\inf\ \left\{\varepsilon\in[0,1]\mid s\sim_{\varepsilon}t\right\}.\end{array}\]\quad  \quad  
  
% END Thmtheorem4.p1
 
% END Thmtheorem4
 
% BEGIN S2.SS2.p5
 
The function $d^{*}$ is a pseudometric on the states of the considered LMP: the triangle inequality comes from the well-known fact, see e.g. [26], that $s\sim_{\varepsilon_{1}}u$ and $u\sim_{\varepsilon_{2}}t$ imply $s\sim_{\varepsilon_{1}+\varepsilon_{2}}t$. (The same property holds for $\prec_{\varepsilon}$ as well.) Indeed, $\prec_{\varepsilon}$ and $\sim_{\varepsilon}$ are not transitive. Instead, these relations are entourages that form a uniform structure, or uniformity, cf. [15].\par
 
% END S2.SS2.p5
 
% BEGIN S2.SS2.p6
 
The kernel of $d^{*}$, that is, the set of pairs at zero distance, is $\sim_{0}$, the bisimilarity relation on LMPs.\par
 
% END S2.SS2.p6
 
% BEGIN Thmtheorem5
 
\par\medskip\noindent\textbf{Theorem 5 ([26]).}\par
 
% BEGIN Thmtheorem5.p1
 
$d^{*}(s,t)=0\Leftrightarrow s\sim_{0}t$.\par
 
% END Thmtheorem5.p1
 
% END Thmtheorem5
 
% BEGIN S2.SS2.p7
 
One of the nice properties of this distance is its relatively easy computability, especially regarding computation “by hand”. It suffices to define an $\varepsilon$-bisimulation to obtain an upper bound to the distance between states. In Example 3, we have $d^{*}(s,t)=\gamma$, and, indeed, we easily see that the probabilities of $s$ and $t$ seen as processes are within $\gamma$. Also here, the Kantorovich distance between these states is 1. One way to see that is by using the functional $\langle a\rangle^{n}1$, which evaluates to $(1-\varepsilon)^{n}$ on $s$, and to 1 on $t$. So $d_{K}(s,t)=1$. Similarly, as already noted in the introduction, for the channel in Figure 1, we can prove that $d^{*}(c_{\varepsilon},c_{\gamma})=|\varepsilon-\gamma|=d_{K}(c_{\varepsilon},c_{\gamma})$, for $\varepsilon,\gamma>0$, but $d^{*}(c_{\varepsilon},c_{0})=\varepsilon\neq d_{K}(c_{\varepsilon},c_{0})$. In fact, $\varepsilon$-distance is incomparable to both the discounted and the undiscounted Kantorovich distance, as argued in [26]. That paper also gives a straightforward polynomial algorithm to compute $d^{*}$. The computation of the Kantorovich distance, discounted or not, has been the subject of many papers, which also led to polynomial algorithms [53].\par
 
% END S2.SS2.p7
 
% END S2.SS2
 
% END S2


% BEGIN S3
 
\par\medskip\noindent\textbf{3 The Lévy-Prokhorov distance lifting}\par
 
% BEGIN S3.p1
 
We will now explain in details what we think should be a behavioral pseudometric on LMPs. These details may not be new, or may appear obvious to the expert reader, but we find it useful to spell them out here. When looking for a distance $d^{\dagger}$ on the states of an LMP, one observes that states are both targets of distributions and (a set of) distributions themselves (since they are defined by their outgoing transition distributions): that is, in Example 3, the states $s$ and $t$ can be viewed as distributions over the set $\{s_{1},s_{2},t_{1}\}$. As there are already a few distances on distributions that were studied outside computer science and concurrency theory, they give a starting point for distances on states: we could say, given such a distance $\delta$:\par
  
\par  \quad  \[d^{\dagger}(s,t):=\sup_{a\in\mathcal{A}}\,\delta(\tau_{a}(s),\tau_{a}(t)),\qquad\mbox{ for $s,t\in S$.}\]\quad  \quad  (2)\quad  
  
We will use the letter $d$ for distances on states, and the symbol $\delta$ for distances on distributions. Analysing the options of this equation in full generality is outside the scope of the current paper but some distances, such as $\delta_{K}^{d}$ (and the one we introduce below $\delta^{d}_{LP}$) have a particularity. They are defined using a parameter $d$, a basic distance on the space where the distributions are defined. This is desirable for concurrency theory because distributions on states that are different but close (or even bisimilar!) should also be close: so a starting distance $d$ on the state space is of key importance to account for these similarities between states/processes. Note that a few papers [30, 19, 51] start from distances on the states that account for extra information (e.g. a distance on the labels/observations that are attached to states) – but this goes beyond the behaviour of states that we want to capture here (although they constitute an interesting line of work to extend our method).\par
 
% END S3.p1
 
% BEGIN S3.p2
 
A needed property for $d^{\dagger}$ to be a behavioural distance, as also pointed out in [9] for $\delta_{K}^{d}$, is that it accords with the starting distance $d$ used in $\delta^{d}$, as follows.\par
  
\par  \quad  \[d^{\dagger}(s,t)=\sup_{a\in\mathcal{A}}\,\delta^{d^{\dagger}}(\tau_{a}(s),\tau_{a}(t)),\qquad\mbox{ for $s,t\in S$.}\]\quad  \quad  (3)\quad  
  
This says that the distance between states viewed as simple members of the space (the left hand part of the equality), is the same as their distance when viewed as processes, that is, their outgoing transition distributions. So this really says that $d^{\dagger}$ treats states according to their behaviour. Technically, this is saying that we are looking for a distance fixpoint $d^{\dagger}$ of the functional\par
  
\par  \quad  \[\Delta(d)(s,t):=\sup_{a\in\mathcal{A}}\,\delta^{d}(\tau_{a}(s),\tau_{a}(t)),\qquad\mbox{ for $s,t\in S$.}\]\quad  \quad  (4)\quad  
  
Of course, by taking $\delta^{d}:=\delta_{K}^{d}$, this instantiates exactly to Equation (1), the fixpoint property of the Kantorovich behavioural distance, observed by [25, 59] and explicitly proven in [57]. One of our goals in this work was to find a similar functional that would have the $\varepsilon$-distance $d^{*}$ as its greatest fixed point. That is, we are seeking a suitable distance $\delta^{d}$ on distributions.\par
 
% END S3.p2
 
% BEGIN Thmtheorem6
 
\par\medskip\noindent\textbf{$\blacktriangleright$ Remark 6.}\par
 
% BEGIN Thmtheorem6.p1
 
Another property that one might consider natural and one might expect from such a distance $\delta^{d}$ is that the distance between the Dirac distributions on states is the same as the starting distance between them. That is, for a distance $d$ on states, one might expect\par
  
\par  \quad  \[d(s,t)=\delta^{d}(1_{s},1_{t}),\qquad\mbox{ for $s,t\in S$,}\]\quad  \quad  (5)\quad  
  
where we write $1_{s}$ for the Dirac distribution on $s\in S$. This is an interesting property of the lifting $\delta^{d}$ – it is a stronger version of the non-expansiveness of the unit of the distribution monad, showing that the unit is an isometry, and we will return to it in Section 3.4 below where we show that it holds for the Lévy-Prokhorov distance, for any starting distance $d$ on states. It also holds for the (undiscounted) Kantorovich distance – this is easy to prove using the Wasserstein formulation of the Kantorovich distance, whereas it does not hold for the total variation distance [43]. Therefore, this property is neither unique to the Lévy-Prokhorov distance that we are interested in, nor it is “behavioural” in the sense that we explained so far: it does not take into account the transition structure $\tau$ of the states. For a somewhat behavioural explanation, note that combining it with Equation (2), this condition implies that a pair of states $s^{\prime},t^{\prime}$ each having a single outgoing transition as a Dirac on state $s$ and $t$ respectively, would get the same distance as $s$ and $t$. In particular, such a distance does not discount the future.\par
 
% END Thmtheorem6.p1
 
% END Thmtheorem6
 
% BEGIN S3.SS1
 
\par\medskip\noindent\textbf{3.1 The Lévy-Prokhorov distance on distributions}\par
 
% BEGIN S3.SS1.p1
 
While examining equation (3) and trying to make $d^{*}$ fit into it as $d^{\dagger}$ we came up with the following distance on distributions. Only afterwards, we discovered that the distance was actually known as the Lévy-Prokhorov lifting of a distance $d$ to distributions. In this section we introduce the distance and provide an example that illustrates it.\par
 
% END S3.SS1.p1
 
% BEGIN Thmtheorem7
 
\par\medskip\noindent\textbf{Definition 7 (Lévy-Prokhorov distance [45]).}\par
 
% BEGIN Thmtheorem7.p1
 
Let $\langle S,d\rangle$ be a pseudometric space; we endow $\mathcal{D}S$ with the pseudometric (or distance) $\delta^{d}_{LP}:\mathcal{D}S\times\mathcal{D}S\to[0,1]$, defined as\par
  
\par  \quad  \[\delta^{d}_{LP}(\mu_{0},\mu_{1})=\inf\{\varepsilon\mid\forall X\subseteq S:\mu_{i}(X)\leq\mu_{1-i}(X^{d}_{\varepsilon})+\varepsilon,\mbox{ for }i=0,1\},\]\quad  \quad  
  
where $X^{d}_{\varepsilon}=\{y\mid\exists x\in X:d(x,y)<\varepsilon\}$. We call $\delta^{d}_{LP}$ the Lévy-Prokhorov (LP, for short) distance.\par
 
% END Thmtheorem7.p1
 
% END Thmtheorem7
 
% BEGIN S3.SS1.p2
 
The strict inequality in the definition of $X_{\varepsilon}^{d}$ is necessary for the proof of Theorem 11.\par
 
% END S3.SS1.p2
 
% BEGIN S3.SS1.p3
 
In the following example we define simple probability measures that help illustrate the need for the extra “$+\,\varepsilon$” in this definition and the need for the “$\varepsilon$ ball around $X$”. At first sight, $\delta^{d}_{LP}$ looks very much like the total variation distance, but this “$\varepsilon$ ball around $X$” makes it very different.\par
 
% END S3.SS1.p3
 
% BEGIN Thmtheorem8
 
\par\medskip\noindent\textbf{Example 8.}\par
 
% BEGIN Thmtheorem8.p1
 
Consider the set of states $S=\{x_{\gamma}\mid\gamma\in[0,1]\}$ and a distance on these states given by $d(x_{\gamma},x_{\xi})=|\gamma-\xi|$, for $\gamma,\xi\in[0,1]$. In a probabilistic transition system, one could imagine that $x_{\gamma}$ has an $a$-loop to itself with probability $1-\gamma$, $\gamma\in[0,1]$, as depicted below (adding transitions out of the states to help see their differences as processes in a transition system). We now define a family of distributions on these states, and we picture them below as states with an outgoing transition without label (one could imagine an $a$-label). Let $\nu_{\gamma}=\gamma 1_{x_{1}}+(1-\gamma)1_{x_{\gamma}}$, for $\gamma\in[0,\frac{1}{2})$, with $1_{x}$ the Dirac measure on $x$. The distributions $\nu_{\gamma}$ and $\nu_{0}$ are illustrated below. For a fixed $\gamma$, these two distributions are actually non-zero on a three-state space $S^{\prime}=\{x_{0},x_{1},x_{\gamma}\}$ .\par
 
% END Thmtheorem8.p1
 
% BEGIN Thmtheorem8.p2
 
\begin{center}\includegraphics[width=0.95\linewidth,height=0.45\textheight,keepaspectratio]{graphic-03.pdf}\end{center}
 
% END Thmtheorem8.p2
 
% BEGIN Thmtheorem8.p3
 
We show that $\delta^{d}_{LP}(\nu_{\gamma},\nu_{0})=\gamma$. Let $1-\gamma>\varepsilon>\gamma$. We show that all such $\varepsilon$ satisfy the inequalities in the definition of $\delta^{d}_{LP}$ and any $\varepsilon\leq\gamma$ does not, so the distance, being the infimum over all such values, is indeed $\gamma$. In particular, we need to check that for all $X\subseteq S$, $\nu_{\gamma}(X)\leq\nu_{0}(X_{\varepsilon}^{d})+\varepsilon$ and $\nu_{0}(X)\leq\nu_{\gamma}(X_{\varepsilon}^{d})+\varepsilon$. The cases $X=S$ and $X=\emptyset$ are trivial, since the distributions are full. For the sets $X$ of size two, we have      \[\begin{array}[]{lllll}\mbox{For }X:=\{x_{0},x_{1}\}:&\nu_{\gamma}(X)=\gamma&\!\!\leq 1+\varepsilon&=\nu_{0}(S)+\varepsilon&=\nu_{0}(X^{d}_{\varepsilon})+\varepsilon\\
&\nu_{0}(X)=1&\!\!\leq 1+\varepsilon&=\nu_{\gamma}(S)+\varepsilon&=\nu_{\gamma}(X^{d}_{\varepsilon})+\varepsilon\\
\mbox{For }X:=\{x_{1},x_{\gamma}\}:&\nu_{\gamma}(X)=1&\!\!\leq 1+\varepsilon&=\nu_{0}(S)+\varepsilon&=\nu_{0}(X^{d}_{\varepsilon})+\varepsilon\\
&\nu_{0}(X)=0&\!\!\leq 1+\varepsilon&=\nu_{\gamma}(S)+\varepsilon&=\nu_{\gamma}(X^{d}_{\varepsilon})+\varepsilon\\
\mbox{For }X:=\{x_{0},x_{\gamma}\}:&\nu_{\gamma}(X)=1-\gamma&\!\!\leq 1+\varepsilon&=\nu_{0}(\{x_{0},x_{\gamma}\})+\varepsilon&=\nu_{0}(X^{d}_{\varepsilon})+\varepsilon\\
&\nu_{0}(X)=1&\!\!\leq(1-\gamma)+\varepsilon&=\nu_{\gamma}(\{x_{0},x_{\gamma}\})+\varepsilon&=\nu_{\gamma}(X^{d}_{\varepsilon})+\varepsilon\\
\end{array}\]\par
 
% END Thmtheorem8.p3
 
% BEGIN Thmtheorem8.p4
 
The last inequality is satisfied thanks to the room given by the “$+\varepsilon$”. For sets of size one, omitting the checks where the probability is zero on the left-hand side of the inequality:\par
  
\par  \quad  \[\begin{array}[]{lllll}\mbox{For }X:=\{x_{1}\}:&\nu_{\gamma}(X)=\gamma&\!\!\leq
0+\varepsilon&=\nu_{0}(\{x_{1}\})+\varepsilon&=\nu_{0}(X^{d}_{\varepsilon})+\varepsilon\\
\mbox{For }X:=\{x_{\gamma}\}:&\nu_{\gamma}(X)=1-\gamma&\!\!\leq 1+\varepsilon&=\nu_{0}(\{x_{\gamma},x_{0}\})+\varepsilon&=\nu_{0}(X^{d}_{\varepsilon})+\varepsilon\\
\mbox{For }X:=\{x_{0}\}:&\nu_{0}(X)=1&\!\!\leq 1-\gamma+\varepsilon&=\nu_{\gamma}(\{x_{\gamma},x_{0}\})+\varepsilon&=\nu_{\gamma}(X^{d}_{\varepsilon})+\varepsilon.\end{array}\]\quad  \quad  
  
In the one but last inequality, it is crucial that $x_{\gamma}$ be within $\varepsilon$ of $x_{0}$, which includes $x_{0}$ in the $\varepsilon$ ball around $x_{\gamma}$. With $X$ in place of $X^{d}_{\varepsilon}$, the inequality would not be satisfied, and not even with $X^{d}_{\gamma/2}$, a set where states at distance zero are included (like bisimilar states – of which there are none here), because we would obtain:\par
  
\par  \quad  \[\nu_{\gamma}(\{x_{\gamma}\})=1-\gamma\not\leq 0+\varepsilon=\nu_{0}(\{x_{\gamma}\}^{d}_{\gamma/2})+\varepsilon.\]\quad  \quad  
  
The last inequality also illustrates this, and also the need for the “$+\varepsilon$”, as the $\varepsilon$-ball itself is not enough for the inequality to be satisfied. Moreover, for a pair of distributions $\nu_{\gamma},\nu_{\xi}$ in our family one can similarly prove that $\delta^{d}_{LP}(\nu_{\gamma},\nu_{\xi})=|\gamma-\xi|$.\par
 
% END Thmtheorem8.p4
 
% BEGIN Thmtheorem8.p5
 
Finally, we note that for $\varepsilon\leq\gamma$ we have $X_{\varepsilon}^{d}=X$ for any $X\subseteq S$ and hence, e.g., for $X=\{x_{\gamma}\}$, as already noted above we have $\nu_{\gamma}(X)=1-\gamma\not\leq 0+\varepsilon=\nu_{0}(X^{d}_{\varepsilon})+\varepsilon$ as $\varepsilon\leq\gamma<\frac{1}{2}$.\par
 
% END Thmtheorem8.p5
 
% END Thmtheorem8
 
% BEGIN S3.SS1.p4
 
One can notice that when viewed as processes, $\nu_{\gamma}$ and $\nu_{0}$ are the same as $s$ and $t$ of Example 3. Technically, it is rather $\tau_{a}(s)$ and $\tau_{a}(t)$ of course.\par
 
% END S3.SS1.p4
 
% BEGIN S3.SS1.p5
 
Once one has seen the Lévy-Prokhorov distance on distributions it seems not surprising that it has some link with the $\varepsilon$-distance. However, our surprise was the other way around, when we realized that the distance on distributions $\delta^{d}$ that we had to concoct to make $d^{*}$ a fixed point of $\Delta$ in Equation (4) already existed.\par
 
% END S3.SS1.p5
 
% END S3.SS1
 
% BEGIN S3.SS2
 
\par\medskip\noindent\textbf{3.2 The Lévy-Prokhorov distance on LMPs}\par
 
% BEGIN S3.SS2.p1
 
We are now ready to define the functional on distances over LMPs that we were looking for. We expect a fixpoint of this functional to be a behavioral distance, but more importantly we expect $d^{*}$ to be the greatest fixpoint of it.\par
 
% END S3.SS2.p1
 
% BEGIN Thmtheorem9
 
\par\medskip\noindent\textbf{Definition 9.}\par
 
% BEGIN Thmtheorem9.p1
 
Let $\langle S,d\rangle$ be a pseudometric space, on which an LMP $\tau:S\to(\mathcal{D}S)^{\mathcal{A}}$ is defined. Let $s_{0}$ and $s_{1}$ be states. We define the functional $\Delta_{LP}$ as\par
  
\par  \quad  $\displaystyle\Delta_{LP}(d)(s_{0},s_{1})$\quad  $\displaystyle=\sup_{a\in\mathcal{A}}\delta_{LP}^{d}(\tau_{a}(s_{0}),\tau_{a}(s_{1}))$\quad  \quad  
 
\par  \quad  \quad  $\displaystyle=\sup_{a\in\mathcal{A}}\inf\{\varepsilon\mid(\forall A\subseteq S:\tau_{a}(s_{i})(A)\leq\tau_{a}(s_{1-i})(A^{d}_{\varepsilon})+\varepsilon,\quad i=0,1)\}.$\quad  \quad  
  
% END Thmtheorem9.p1
 
% END Thmtheorem9
 
% BEGIN S3.SS2.p2
 
Hence, for a fixed pseudometric space $\langle S,d\rangle$, $\Delta_{LP}$ applied to $d$ gives a new pseudometric on $S$. Clearly, this can be seen as a functional on $\mathbf{PMet_{1}}$, mapping $\langle S,d\rangle$ to $(S,\Delta_{LP}(d))$.\par
 
% END S3.SS2.p2
 
% BEGIN S3.SS2.p3
 
Before proving that $d^{*}$ is a fixpoint, we show how the supremum over actions can safely be replaced by inserting a universal quantification on actions inside the set, as will be useful in Proposition 15.\par
 
% END S3.SS2.p3
 
% BEGIN Thmtheorem10
 
\par\medskip\noindent\textbf{Lemma 10.}\par
 
% BEGIN Thmtheorem10.p1
 
Let $\langle S,d\rangle$ be a pseudometric space, on which an LMP $\tau:S\to(\mathcal{D}S)^{\mathcal{A}}$ is defined. Let $s_{0}$ and $s_{1}$ be states. Then\par
  
\par  \quad  $\displaystyle\Delta_{LP}(d)(s_{0},s_{1})$\quad  $\displaystyle=\inf\{\varepsilon\mid(\forall a\in\mathcal{A},\forall A\subseteq
S:\tau_{a}(s_{i})(A)\leq\tau_{a}(s_{1-i})(A^{d}_{\varepsilon})+\varepsilon,\quad i=0,1)\}.$\quad  \quad  
  
% END Thmtheorem10.p1
 
% END Thmtheorem10
 
% BEGIN S3.SS2.p4
 
One remarkable property of prefixpoints of this functional is that they define $\varepsilon$-bisimulations. Of course the theorem is also true for fixpoints.\par
 
% END S3.SS2.p4
 
% BEGIN Thmtheorem11
 
\par\medskip\noindent\textbf{Theorem 11.}\par
 
% BEGIN Thmtheorem11.p1
 
Let $\tau:S\to(\mathcal{D}S)^{\mathcal{A}}$ be an LMP with $\langle S,d\rangle$ a pseudometric space. If $d$ is a prefixpoint of $\Delta_{LP}$, then for any $\varepsilon>0$, $R_{\varepsilon}:=\{(s,t)\in S\times S\mid d(s,t)<\varepsilon\}$, is an $\varepsilon$-bisimulation.\par
 
% END Thmtheorem11.p1
 
% END Thmtheorem11
 
% BEGIN S3.SS2.SSS0.P0.SPx1
 
\par\medskip\noindent\textbf{Proof.}\par
 
% BEGIN S3.SS2.SSS0.P0.SPx1.p1
 
Let $d$ be a prefixpoint, i.e., let $\Delta_{LP}(d)\leq d$. Let $R_{\varepsilon}=\{(s,t)\mid d(s,t)<\varepsilon\}$. We show it is an $\varepsilon$-bisimulation. Let $s_{0}R_{\varepsilon}s_{1}$, i.e., $d(s_{0},s_{1})<\varepsilon$, and also $\Delta_{LP}(d)(s_{0},s_{1})<\varepsilon$, because $d$ is a prefixpoint. Let $a\in\mathcal{A}$; then $\delta_{LP}^{d}(\tau_{a}(s_{0}),\tau_{a}(s_{1}))<\varepsilon$. Because $\delta_{LP}^{d}$ is an infimum, there is some $\gamma<\varepsilon$ such that for $i=1,2$ and $X\subseteq S$\par
  
\par  \quad  $\displaystyle\tau_{a}(s_{i})(X)$\quad  $\displaystyle\leq\tau_{a}(s_{1-i})(X^{d}_{\gamma})+\gamma$\quad  by definition of $\delta_{LP}$\quad  \quad  
 
\par  \quad  \quad  $\displaystyle\leq\tau_{a}(s_{1-i})(X^{d}_{\varepsilon})+\varepsilon$\quad  $\displaystyle\mbox{because }X^{d}_{\gamma}\subseteq X^{d}_{\varepsilon}\mbox{ and }\gamma\leq\varepsilon.$\quad  \quad  
 
\par  \quad  \quad  $\displaystyle=\tau_{a}(s_{1-i})(R_{\varepsilon}(X))+\varepsilon$\quad  $\displaystyle\mbox{by definition of }R_{\varepsilon}(X).$\quad  \quad  
  
So $R_{\varepsilon}$ is an $\varepsilon$-bisimulation, as wanted. $\hfill\blacktriangleleft$ The proof of this theorem needs that $X_{\varepsilon}^{d}$ be defined with a strict inequality. It is necessary for the existence of $\gamma$ (and for the last equality).\par
 
% END S3.SS2.SSS0.P0.SPx1.p1
 
% BEGIN S3.SS2.SSS0.P0.SPx1.p2
 
This theorem is nice in itself but it also has an important corollary, that $d^{*}$ is greater than or equal to all fixpoints of $\Delta_{LP}$.\par
 
% END S3.SS2.SSS0.P0.SPx1.p2
 
% BEGIN Thmtheorem12
 
\par\medskip\noindent\textbf{Corollary 12.}\par
 
% BEGIN Thmtheorem12.p1
 
$d^{*}$ is greater than or equal to all fixpoints, i.e.,\par
  
\par  \quad  \[\Delta_{LP}(d)=d\quad\Longrightarrow\quad d\sqsubseteq d^{*}.\]\quad  \quad  
  
% END Thmtheorem12.p1
 
% END Thmtheorem12
 
% END S3.SS2.SSS0.P0.SPx1
 
% BEGIN S3.SS2.SSS0.P0.SPx2
 
\par\medskip\noindent\textbf{Proof.}\par
 
% BEGIN S3.SS2.SSS0.P0.SPx2.p1
 
Let $d$ be a fixpoint of $\Delta_{LP}$ and let $d(s,t)=\varepsilon$. By the previous theorem, $s\sim_{\gamma}t$ for all $\gamma>\varepsilon$. Since $d^{*}(s,t)=\inf\{e\mid s\sim_{e}t\}$, we obtain $d^{*}(s,t)\leq\varepsilon=d(s,t)$, as wanted. $\hfill\blacktriangleleft$\par
 
% END S3.SS2.SSS0.P0.SPx2.p1
 
% BEGIN S3.SS2.SSS0.P0.SPx2.p2
 
With the next proposition, we provide another example of a fixpoint of $\Delta_{LP}$: the exact bisimilarity distance, which we denote by $d_{\sim}$. It is equal to 0 if the states are bisimilar, otherwise it is 1.\par
 
% END S3.SS2.SSS0.P0.SPx2.p2
 
% BEGIN Thmtheorem13
 
\par\medskip\noindent\textbf{Proposition 13.}\par
 
% BEGIN Thmtheorem13.p1
 
The exact bisimilarity distance $d_{\sim}$ is a fixpoint of $\Delta_{LP}$.\par
 
% END Thmtheorem13.p1
 
% END Thmtheorem13
 
% END S3.SS2.SSS0.P0.SPx2
 
% BEGIN S3.SS2.SSS0.P0.SPx3
 
\par\medskip\noindent\textbf{Proof.}\par
 
% BEGIN S3.SS2.SSS0.P0.SPx3.p1
 
We first prove that bisimilar states $s,t$ satisfy $\Delta_{LP}(d_{\sim})(s,t)=0$. Consider $R$ to be bisimilarity. Because $R$ is a bisimulation, we have that for all $s_{0}Rs_{1}$, $a\in\mathcal{A}$ and $X\subseteq S$,\par
  
\par  \quad  \[\tau_{a}(s_{0})(X)\leq\tau_{a}(s_{1})(R(X))=\tau_{a}(s_{1})(X^{d_{\sim}}_{\gamma}),\]\quad  \quad  
  
for any $\gamma\in(0,1)$, because $X^{d_{\sim}}_{\gamma}$ is just the smallest $R$-closed set that contains $X$. A symmetric argument applies with 0 and 1 interchanged. Consequently, $\Delta_{LP}(d_{\sim})(s_{0},s_{1})=0$, as wanted.\par
 
% END S3.SS2.SSS0.P0.SPx3.p1
 
% BEGIN S3.SS2.SSS0.P0.SPx3.p2
 
Conversely, if $\Delta_{LP}(d_{\sim})(s_{0},s_{1})=0$, then there is a sequence of $\gamma_{n}\in(0,1)$ that converges to zero, for which we have\par
  
\par  \quad  \[\tau_{a}(s_{0})(X)\leq\tau_{a}(s_{1})(X^{d_{\sim}}_{\gamma_{n}})+\gamma_{n}=\tau_{a}(s_{1})(R(X))+\gamma_{n}.\]\quad  \quad  
  
Since the property is monotone, recall also (6), for all $\gamma\in(0,1)$ we have\par
  
\par  \quad  \[\tau_{a}(s_{0})(X)\leq\tau_{a}(s_{1})(X^{d_{\sim}}_{\gamma})+\gamma=\tau_{a}(s_{1})(R(X))+\gamma.\]\quad  \quad  
  
So $\tau_{a}(s_{0})(X)\leq\tau_{a}(s_{1})(R(X))$, as wanted. $\hfill\blacktriangleleft$\par
 
% END S3.SS2.SSS0.P0.SPx3.p2
 
% BEGIN Thmtheorem14
 
\par\medskip\noindent\textbf{Proposition 14.}\par
 
% BEGIN Thmtheorem14.p1
 
$\Delta_{LP}$ is monotonic.\par
 
% END Thmtheorem14.p1
 
% END Thmtheorem14
 
% BEGIN S3.SS2.SSS0.P0.SPx3.p3
 
We can now prove that the distance defined by $\varepsilon$-bisimulation is a fixpoint of $\Delta_{LP}$.\par
 
% END S3.SS2.SSS0.P0.SPx3.p3
 
% BEGIN Thmtheorem15
 
\par\medskip\noindent\textbf{Proposition 15.}\par
 
% BEGIN Thmtheorem15.p1
 
$d^{*}$ is a fixpoint of $\Delta_{LP}$, and hence it is the greatest fixpoint.\par
 
% END Thmtheorem15.p1
 
% END Thmtheorem15
 
% BEGIN S3.SS2.SSS0.P0.SPx3.p4
 
At this point we remark again that even if it seems not surprising that the Lévy-Prokhorov metric on distributions has some link with the $\varepsilon$-distance, going from the basic metric on distributions $\delta_{LP}$ to a meaningful distance on LMPs is not direct (one has to find the fixed-point properties).\par
 
% END S3.SS2.SSS0.P0.SPx3.p4
 
% END S3.SS2.SSS0.P0.SPx3
 
% END S3.SS2
 
% BEGIN S3.SS3
 
\par\medskip\noindent\textbf{3.3 The Lévy-Prokhorov lifting of the subdistribution functor to $\mathbf{PMet_{1}}$}\par
 
% BEGIN S3.SS3.p1
 
Using the LP distance lifting, we define a functor $\mathcal{D}_{LP}$ on $\mathbf{PMet_{1}}$, the category of 1-bounded pseudometric spaces with nonexpansive functions, as follows: On objects $\langle X,d\rangle$, we have\par
  
\par  \quad  \[\mathcal{D}_{LP}\langle X,d\rangle=(\mathcal{D}X,\delta^{d}_{LP})\]\quad  \quad  
  
and on morphisms $f\colon\langle X_{0},d_{0}\rangle\to\langle X_{1},d_{1}\rangle$ we set $\mathcal{D}_{LP}f=\mathcal{D}f$, that is\par
  
\par  \quad  \[\mathcal{D}_{LP}f\colon\mathcal{D}_{LP}\langle X_{0},d_{0}\rangle\to\mathcal{D}_{LP}\langle X_{1},d_{1}\rangle\textrm{ with }\varphi\mapsto\lambda y.\,\varphi(f^{-1}(\{y\})).\]\quad  \quad  
  
% END S3.SS3.p1
 
% BEGIN Thmtheorem16
 
\par\medskip\noindent\textbf{Proposition 16.}\par
 
% BEGIN Thmtheorem16.p1
 
$\mathcal{D}_{LP}$ is a functor on $\mathbf{PMet_{1}}$.\par
 
% END Thmtheorem16.p1
 
% END Thmtheorem16
 
% BEGIN S3.SS3.p2
 
Just like for the Kantorovich-lifting of the distribution functor, we can prove that the Lévy-Prokhorov lifting of the distribution functor is locally nonexpansive, in the next proposition. For this reason, note that given pseudometric spaces $\langle X,d_{X}\rangle$ and $\langle Y,d_{Y}\rangle$ in $\mathbf{PMet_{1}}$, the hom-set $X\to Y$ of all nonexpansive maps from $X$ to $Y$ carries a metric defined by\par
  
\par  \quad  \[d_{X\to Y}(f_{1},f_{2})=\sup_{x\in X}d_{Y}(f_{1}(x),f_{2}(x)).\]\quad  \quad  
  
% END S3.SS3.p2
 
% BEGIN Thmtheorem17
 
\par\medskip\noindent\textbf{Proposition 17.}\par
 
% BEGIN Thmtheorem17.p1
 
The functor $\mathcal{D}_{LP}$ is locally nonexpansive, that is, for $f_{1},f_{2}\in X\to Y$\par
  
\par  \quad  \[\delta_{\mathcal{D}_{LP}X\to\mathcal{D}_{LP}Y}(\mathcal{D}_{LP}f_{1},\mathcal{D}_{LP}f_{2})\leq d_{X\to Y}(f_{1},f_{2})\]\quad  \quad  
  
with $\delta_{\mathcal{D}_{LP}X\to\mathcal{D}_{LP}Y}$ being the metric on the hom-set $\mathcal{D}_{LP}\langle X,d_{X}\rangle\to\mathcal{D}_{LP}\langle Y,d_{Y}\rangle$.\par
 
% END Thmtheorem17.p1
 
% END Thmtheorem17
 
% BEGIN S3.SS3.p3
 
For the following property, recall that a map $f\colon X\to Y$ for (pseudo)metric spaces $\langle X,d_{X}\rangle$ and $\langle Y,d_{Y}\rangle$ is an isometry if and only if for all $x,y\in X$, $d_{Y}(f(x),f(y))=d_{X}(x,y)$.\par
 
% END S3.SS3.p3
 
% BEGIN Thmtheorem18
 
\par\medskip\noindent\textbf{Proposition 18.}\par
 
% BEGIN Thmtheorem18.p1
 
The functor $\mathcal{D}_{LP}$ preserves isometries, i.e., if $f\colon X\to Y$ is an isometry for (pseudo)metric spaces $\langle X,d_{X}\rangle$ and $\langle Y,d_{Y}\rangle$, then so is $\mathcal{D}_{LP}f$.\par
 
% END Thmtheorem18.p1
 
% END Thmtheorem18
 
% BEGIN S3.SS3.p4
 
As a consequence, from the results of [31], we obtain that the lifting functor $\mathcal{D}_{LP}$ is Kantorovich, i.e., it has codensity lifting which means that there is a set of predicate liftings of type $\mathcal{D}[0,1]\to[0,1]$ such that codensity lifting gives rise to $\mathcal{D}_{LP}$. This has already been remarked in [31, Example 5.5.2]. However, the exact set of predicate liftings needed to give rise to the LP-lifting is not provided.\par
 
% END S3.SS3.p4
 
% END S3.SS3
 
% BEGIN S3.SS4
 
\par\medskip\noindent\textbf{3.4 The lifted functor $\mathcal{D}_{LP}$ is not a monad lifting of $\mathcal{D}$}\par
 
% BEGIN S3.SS4.p1
 
Behavioural distances have been axiomatized within the line of work on quantitative equational theories [39]. It was therefore a natural question for us whether the $\varepsilon$-distance, or the Lévy-Prokhorov distance itself can be given a quantitative axiomatization. This is what we briefly investigate in this section.\par
 
% END S3.SS4.p1
 
% BEGIN S3.SS4.p2
 
Mio et al. [43, Lemma 7.2, Theorem 7.7(2)] have shown that:\par
  
\par\noindent  $\blacksquare$ 
% BEGIN S3.I1.i1.p1
 
If a monad on metric spaces is axiomatizable, then (just by being a monad on metric spaces) it has nonexpansive unit and multiplication.\par
 
% END S3.I1.i1.p1
 
 
\par\noindent  $\blacksquare$ 
% BEGIN S3.I1.i2.p1
 
A monad on metric spaces that is a lifting of a monad on $\mathbf{Sets}$ is axiomatizable with a quantitative theory. Being a lifting means that it acts on objects and arrows in metric spaces in the same way it does in $\mathbf{Sets}$ and the unit and multiplication are nonexpansive.\par
 
% END S3.I1.i2.p1
 
  
The Kantorovich lifting of $\mathcal{D}$ has been axiomatized, and the total variation distance has been shown non-axiomatizable (as the unit is not nonexpansive). In the case of $\mathcal{D}_{LP}$ we can see that the unit is nonexpansive, but the multiplication is not, and hence $\mathcal{D}_{LP}$ is not a lifting of the monad $\mathcal{D}$ to $\mathbf{PMet_{1}}$ and the Lévy-Prokhorov distance on distributions is not axiomatizable with a quantitative theory, at least not with the standard multiplication.\par
 
% END S3.SS4.p2
 
% BEGIN S3.SS4.p3
 
Recall the definitions of $\eta$ and $\mu$ for the subdistribution monad $\mathcal{D}$:\par
  
\par  \quad  \[\eta_{X}(x)=1_{x},\text{ the Dirac distribution at }x;\text{ and }\mu_{X}(\Phi)(x)=\sum_{\varphi}\Phi(\varphi)\cdot\varphi(x).\]\quad  \quad  
  
% END S3.SS4.p3
 
% BEGIN Thmtheorem19
 
\par\medskip\noindent\textbf{Lemma 19.}\par
 
% BEGIN Thmtheorem19.p1
 
The unit $\eta$ of $\mathcal{D}$ is nonexpansive with respect to the Lévy-Prokhorov distance. Moreover, it is an isometry, i.e., $\delta_{LP}^{d}(1_{x},1_{y})=d(x,y)$.\par
 
% END Thmtheorem19.p1
 
% END Thmtheorem19
 
% BEGIN S3.SS4.p4
 
The unit of the distribution monad is also an isometry with respect to the Kantorovich distance (without discount), as is easy to prove using the duality with the Wasserstein distance. However, the multiplication $\mu$ of the monad $\mathcal{D}$ is not nonexpansive, as the following example shows.\par
 
% END S3.SS4.p4
 
% BEGIN Thmtheorem20
 
\par\medskip\noindent\textbf{Example 20.}\par
 
% BEGIN Thmtheorem20.p1
 
Consider $a$, $b$ and $c$ in the figure below as distributions over $\{\bot,\bullet\}$.\par
 
% END Thmtheorem20.p1
 
% BEGIN Thmtheorem20.p2
 
\begin{center}\includegraphics[width=0.95\linewidth,height=0.45\textheight,keepaspectratio]{graphic-04.pdf}\end{center}
 
% END Thmtheorem20.p2
 
% BEGIN Thmtheorem20.p3
 
The ambient distance is $d(\bot,\bullet)=1$, we omit it from the notation: we write $\delta_{LP}$ for $\delta_{LP}^{d}$. Consider the distributions $\varphi:=1_{a}$ and $\psi:=(1-\varepsilon)1_{b}+\varepsilon 1_{c}$ over these distributions, as pictured. We collect several useful facts:\par
 
% END Thmtheorem20.p3
 
% BEGIN Thmtheorem20.p4
  
\par\noindent  $\blacksquare$ 
% BEGIN S3.I2.i1.p1
 
$\delta_{LP}(a,b)=\varepsilon$ and $\delta_{LP}^{\delta_{LP}}(\varphi,\psi)=\varepsilon$.\par
 
% END S3.I2.i1.p1
 
 
\par\noindent  $\blacksquare$ 
% BEGIN S3.I2.i2.p1
 
$\mu\varphi=a$ and $\mu\psi=(1-\varepsilon)^{2}\bot+(1-\varepsilon)\varepsilon\bullet+\varepsilon\bullet=(1-\varepsilon)^{2}\bot+(2-\varepsilon)\varepsilon\bullet$.\par
 
% END S3.I2.i2.p1
 
 
\par\noindent  $\blacksquare$ 
% BEGIN S3.I2.i3.p1
 
$\{\bot\}_{\gamma}=\{\bot\}$ for any $\gamma\leq\varepsilon$ and hence $\mu\varphi(\{\bot\})=1$ and $\mu\psi(\{\bot\}_{\gamma})=\mu\psi(\{\bot\})=(1-\varepsilon)^{2}$.\par
 
% END S3.I2.i3.p1
 
  
% END Thmtheorem20.p4
 
% BEGIN Thmtheorem20.p5
 
For $\gamma<\varepsilon(2-\varepsilon)$, we have $1\not\leq(1-\varepsilon)^{2}+\gamma$ which therefore yields\par
  
\par  \quad  \[\mu\varphi(\{\bot\})=1\not\leq(1-\varepsilon)^{2}+\gamma=\mu\psi(\{\bot\}_{\gamma})+\gamma.\]\quad  \quad  
  
This shows that $\delta_{LP}(\mu\varphi,\mu\psi)\geq\varepsilon(2-\varepsilon)>\varepsilon=\delta_{LP}^{\delta_{LP}}(\varphi,\psi)$ and hence $\mu$ is not nonexpansive.\par
 
% END Thmtheorem20.p5
 
% END Thmtheorem20
 
% BEGIN S3.SS4.p5
 
This result is not surprising, as the binding operator $\mu$ really is a multiplication, and so it accords well with the Kantorovich metric, which does multiply the probabilities of distributions in play. We do not know yet if another operator could help express the functor $\mathcal{D}_{LP}$ as a monad lifting.\par
 
% END S3.SS4.p5
 
% END S3.SS4
 
% END S3


% BEGIN A1.SS0.SSS0.P0.SPx1
 
\par\medskip\noindent\textbf{Proof.}\par
 
% BEGIN A1.SS0.SSS0.P0.SPx1.p1
 
For $a\in\mathcal{A}$ and $\varepsilon>0$, we define the following predicate:\par
  
\par  \quad  \[P(a,\varepsilon):=(\forall A\subseteq S:\tau_{a}(s_{i})(A)\leq\tau_{a}(s_{1-i})(A^{d}_{\varepsilon})+\varepsilon,\quad i=0,1).\]\quad  \quad  
  
This family of predicates satisfies\par
  
\par  \quad  \[\mbox{if }\gamma>\varepsilon,P(a,\varepsilon)\Rightarrow P(a,\gamma).\]\quad  \quad  (6)\quad  
  
We now prove that\par
  
\par  \quad  \[\sup_{a\in\mathcal{A}}\inf\{\varepsilon\mid P(a,\varepsilon)\}=\inf\{\varepsilon\mid(\forall a\in\mathcal{A}:P(a,\varepsilon))\}\]\quad  \quad  
   $\leq$:  
% BEGIN A1.I1.ix1.p1
 
The following sequence of arguments shows this inequality.\par
  
\par  \quad  $\displaystyle\{\varepsilon\mid P(a_{0},\varepsilon)\}$\quad  $\displaystyle\supseteq\{\varepsilon\mid(\forall a\in\mathcal{A}:P(a,\varepsilon))\}$\quad  $\displaystyle\mbox{for all }a_{0}\in\mathcal{A},\mbox{ so}$\quad  \quad  
 
\par  \quad  $\displaystyle\inf\{\varepsilon\mid P(a_{0},\varepsilon)\}$\quad  $\displaystyle\leq\inf\{\varepsilon\mid(\forall a\in\mathcal{A}:P(a,\varepsilon))\}$\quad  $\displaystyle\mbox{for all }a_{0}\in\mathcal{A},\mbox{ and so}$\quad  \quad  
 
\par  \quad  $\displaystyle\sup_{a\in\mathcal{A}}\inf\{\varepsilon\mid P(a,\varepsilon)\}$\quad  $\displaystyle\leq\inf\{\varepsilon\mid(\forall a\in\mathcal{A}:P(a,\varepsilon))\}.$\quad  \quad  
  
% END A1.I1.ix1.p1
  $\geq$:  
% BEGIN A1.I1.ix2.p1
 
Let $\alpha=\sup_{a\in\mathcal{A}}\inf\{\varepsilon\mid P(a,\varepsilon)\}.$ Then for all $a\in\mathcal{A}$, $\alpha\geq\inf\{\varepsilon\mid P(a,\varepsilon)\}.$ Let $\gamma>\alpha$. Then we have\par
  
\par  \quad  $\displaystyle\mbox{for all }a\in\mathcal{A}$\quad  $\displaystyle\quad\gamma\in\{\varepsilon\mid P(a,\varepsilon)\}$\quad  $\displaystyle\mbox{by\nobreak\ (6) and because }\gamma>\inf\{\varepsilon\mid P(a,\varepsilon)\}$\quad  \quad  
 
\par  \quad  so\quad  $\displaystyle\quad\gamma\in\cap_{a\in\mathcal{A}}\{\varepsilon\mid P(a,\varepsilon)\}$\quad  \quad  
 
\par  \quad  so\quad  $\displaystyle\quad\gamma\in\{\varepsilon\mid(\forall a\in\mathcal{A}:P(a,\varepsilon))\}$\quad  \quad  
 
\par  \quad  so\quad  $\displaystyle\quad\gamma\geq\inf\{\varepsilon\mid(\forall a\in\mathcal{A}:P(a,\varepsilon))\}.\hskip-11.38109pt$\quad  \quad  
  
This implies $\alpha\geq\inf\{\varepsilon\mid(\forall a\in\mathcal{A}:P(a,\varepsilon))\}$ as well. Otherwise, suppose $\alpha<\iota:=\inf\{\varepsilon\mid(\forall a\in\mathcal{A}:P(a,\varepsilon))\}$. Then $\alpha<(\alpha+\iota)/2<\iota$, so the result above with $\gamma=(\alpha+\iota)/2$ yields $(\alpha+\iota)/2\geq\iota$, a contradiction.\par
 
% END A1.I1.ix2.p1
   
$\hfill\blacktriangleleft$\par
 
% END A1.SS0.SSS0.P0.SPx1.p1
 
% BEGIN A1.SS0.SSS0.P0.SPx1.p2
 
The proof that $\Delta_{LP}$ is monotonic needs the following simple lemma.\par
 
% END A1.SS0.SSS0.P0.SPx1.p2
 
% BEGIN Thmtheorem29
 
\par\medskip\noindent\textbf{Lemma 29.}\par
 
% BEGIN Thmtheorem29.p1
 
If $d_{2}\sqsubseteq d_{1}$, we have $X^{d_{2}}_{\varepsilon}\subseteq X^{d_{1}}_{\varepsilon}$.\par
 
% END Thmtheorem29.p1
 
% END Thmtheorem29
 
% END A1.SS0.SSS0.P0.SPx1


% BEGIN A1.SS0.SSS0.P0.SPx2
 
\par\medskip\noindent\textbf{Proof.}\par
 
% BEGIN A1.SS0.SSS0.P0.SPx2.p1
 
Assume $d_{1}\leq d_{2}$. Let $y\in X^{d_{2}}_{\varepsilon}$. Then there is some $x\in X$ such that $d_{2}(x,y)<\varepsilon$. This gives $d_{1}(x,y)\leq d_{2}(x,y)<\varepsilon$. So $y\in X^{d_{1}}_{\varepsilon}$. $\hfill\blacktriangleleft$\par
 
% END A1.SS0.SSS0.P0.SPx2.p1
 
% BEGIN A1.SS0.SSS0.P0.SPx2.p2
 
% BEGIN Thmtheorem14
 
\par\medskip\noindent\textbf{Proposition 14. [Restated, see original statement.]}\par
 
% BEGIN Thmtheorem14.p1
 
$\Delta_{LP}$ is monotonic.\par
 
% END Thmtheorem14.p1
 
% END Thmtheorem14
 
% END A1.SS0.SSS0.P0.SPx2.p2
 
% END A1.SS0.SSS0.P0.SPx2


% BEGIN A1.SS0.SSS0.P0.SPx3
 
\par\medskip\noindent\textbf{Proof.}\par
 
% BEGIN A1.SS0.SSS0.P0.SPx3.p1
 
Let $d_{1}\leq d_{2}$. Let $s_{0}$, $s_{1}$ be states, let $a\in\mathcal{A}$ and define the two sets\par
  
\par  \quad  $\displaystyle E_{1}$\quad  $\displaystyle=\{\varepsilon\mid(\forall A\subseteq S:\tau_{a}(s_{i},A)\leq\tau_{a}(s_{1-i},A^{d_{1}}_{\varepsilon})+\varepsilon,\quad i=0,1)\}$\quad  \quad  
 
\par  \quad  $\displaystyle\mbox{and}\quad E_{2}$\quad  $\displaystyle=\{\varepsilon\mid(\forall A\subseteq S:\tau_{a}(s_{i},A)\leq\tau_{a}(s_{1-i},A^{d_{2}}_{\varepsilon})+\varepsilon,\quad i=0,1)\}.$\quad  \quad  
  
Then $E_{2}\subseteq E_{1}$. Indeed, let $\varepsilon\in E_{2}$ and $A\subseteq S$. Then $\tau_{a}(s_{i},A)\leq\tau_{a}(s_{1-i},A^{d_{2}}_{\varepsilon})+\varepsilon\leq\tau_{a}(s_{1-i},A^{d_{1}}_{\varepsilon})+\varepsilon$, because $A^{d_{2}}_{\varepsilon}\subseteq A^{d_{1}}_{\varepsilon}$ by Lemma 29. So $\inf E_{1}\leq\inf E_{2}$ and hence $\Delta_{LP}(d_{1})=\sup_{a\in\mathcal{A}}\inf E_{1}\leq\sup_{a\in\mathcal{A}}\inf E_{2}=\Delta_{LP}(d_{2})$, as wanted. $\hfill\blacktriangleleft$\par
 
% END A1.SS0.SSS0.P0.SPx3.p1
 
% BEGIN A1.SS0.SSS0.P0.SPx3.p2
 
% BEGIN Thmtheorem15
 
\par\medskip\noindent\textbf{Proposition 15. [Restated, see original statement.]}\par
 
% BEGIN Thmtheorem15.p1
 
$d^{*}$ is a fixpoint of $\Delta_{LP}$, and hence it is the greatest fixpoint.\par
 
% END Thmtheorem15.p1
 
% END Thmtheorem15
 
% END A1.SS0.SSS0.P0.SPx3.p2
 
% END A1.SS0.SSS0.P0.SPx3


% BEGIN A1.SS0.SSS0.P0.SPx4
 
\par\medskip\noindent\textbf{Proof.}\par
 
% BEGIN A1.SS0.SSS0.P0.SPx4.p1
 
Let $s_{0}$, $s_{1}\in S$. We want $\Delta_{LP}(d^{*})(s_{0},s_{1})=d^{*}(s_{0},s_{1})$. Consider the following sets\par
  
\par  \quad  $\displaystyle E^{\Delta}$\quad  $\displaystyle=\{\varepsilon\mid\forall a\in\mathcal{A}:(\forall X\subseteq S:\tau_{a}(s_{i},X)\leq\tau_{a}(s_{1-i},X^{d^{*}}_{\varepsilon})+\varepsilon,\quad i=0,1)\}$\quad  \quad  
 
\par  \quad  $\displaystyle\mbox{and}\quad E^{*}$\quad  $\displaystyle=\{\varepsilon\mid s_{0}\sim_{\varepsilon}s_{1}\}$\quad  \quad  
 
\par  \quad  \quad  $\displaystyle=\{\varepsilon\mid\exists R\mbox{ an $\varepsilon$-bisimulation s.t. $s_{0}Rs_{1}$, that is, for all $(t_{0},t_{1})\in R$, we have}$\quad  \quad  
 
\par  \quad  \quad  $\displaystyle\mbox{\nobreak\ \nobreak\ \nobreak\ \nobreak\ }\quad\forall a\in\mathcal{A}:(\forall X\subseteq S:\tau_{a}(t_{i},X)\leq\tau_{a}(t_{1-i},R(X))+\varepsilon,\quad i=0,1)\}.$\quad  \quad  
  
Then $\Delta(d^{*})(s_{0},s_{1})=\inf E^{\Delta}$ by Lemma 10 and $d^{*}(s_{0},s_{1})=\inf E^{*}\!$. We prove $\inf E^{\Delta}=\inf E^{*}\!$.\par
  $\geq$:  
% BEGIN A1.I2.ix1.p1
 
We prove that $E^{\Delta}\subseteq E^{*}$. Let $\varepsilon\in E^{\Delta}$. Let $R_{\varepsilon}=\{(x,y)\mid x\sim_{\varepsilon}y\}$, that is, the biggest $\varepsilon$-bisimulation. We need to prove that $s_{0}\sim_{\varepsilon}s_{1}$. For that we prove that $R_{\varepsilon}^{+}:=R_{\varepsilon}\cup\{(s_{0},s_{1})\}$ is an $\varepsilon$-bisimulation. Let $a\in\mathcal{A}$ and $X\subseteq S$. Then $X_{\varepsilon}^{d^{*}}\subseteq R^{+}_{\varepsilon}(X)$. Indeed, let $y\in X_{\varepsilon}^{d^{*}}$ then $d^{*}(y,x)<\varepsilon$ for some $x\in X$. So, because $d^{*}(y,x)=\inf\{e\mid y\sim_{e}x\}$, we have $y\sim_{\gamma}x$ for some $\gamma<\varepsilon$, and hence $y\sim_{\varepsilon}x$ as every $\gamma$-bisimulation is an $\varepsilon$-bisimulation. Hence $y\in R^{+}_{\varepsilon}(X).$ We only need to prove the bisimulation condition for the pair $(s_{0},s_{1})$. Then, for $i=0,1$ we have\par
  
\par  \quad  $\displaystyle\tau_{a}(s_{i},X)$\quad  $\displaystyle\leq\tau_{a}(s_{1-i},X_{\varepsilon}^{d^{*}})+\varepsilon$\quad  $\displaystyle\mbox{by choice of }\varepsilon$\quad  \quad  
 
\par  \quad  \quad  $\displaystyle\leq\tau_{a}(s_{1-i},R^{+}_{\varepsilon}(X))+\varepsilon$\quad  $\displaystyle\mbox{because }X_{\varepsilon}^{d^{*}}\subseteq R^{+}_{\varepsilon}(X).$\quad  \quad  
  
So $s_{0}\sim_{\varepsilon}s_{1}$ and hence $\varepsilon\in E^{*}$.\par
 
% END A1.I2.ix1.p1
  $\leq$:  
% BEGIN A1.I2.ix2.p1
 
Let $\varepsilon\in E^{*}$. We prove that for all $\gamma>\varepsilon$, we have $\gamma\in E^{\Delta}$. This will prove that $\inf E^{\Delta}\leq\varepsilon$. Since $\varepsilon$ is arbitrary, this will give the result: $\inf E^{\Delta}\leq\inf E^{*}$. So let $\gamma>\varepsilon$. Because $\varepsilon\in E^{*}$, we have $s_{0}\sim_{\varepsilon}s_{1}$, and hence there is some $\varepsilon$-bisimulation $R$ such that $s_{0}Rs_{1}$. We want $\gamma\in E^{\Delta}$, so let $a\in\mathcal{A}$ and $X\subseteq S$. We want $\tau_{a}(s_{i},X)\leq\tau_{a}(s_{1-i},X^{d^{*}}_{\gamma})+\gamma,\quad i=0,1$. First observe that $R(X)\subseteq X^{d^{*}}_{\gamma}$. Indeed, if $y\in R(X)$, then there is some $x\in X$ such that $yRx$. Because $R$ is an $\varepsilon$-bisimulation, $d^{*}(y,x)\leq\varepsilon$, so $y\in X^{d^{*}}_{\gamma}$ (we cannot say $y\in X^{d^{*}}_{\varepsilon}$ because these sets are open balls). Then, for $i=0,1$:\par
  
\par  \quad  $\displaystyle\tau_{a}(s_{i},X)$\quad  $\displaystyle\leq\tau_{a}(s_{1-i},R(X))+\varepsilon,$\quad  because $R$ is an $\varepsilon$-bisimulation\quad  \quad  
 
\par  \quad  \quad  $\displaystyle\leq\tau_{a}(s_{1-i},X^{d^{*}}_{\gamma})+\gamma,$\quad  $\displaystyle\mbox{because }R(X)\subseteq X^{d^{*}}_{\gamma}\mbox{and }\varepsilon<\gamma.$\quad  \quad  
  
So, $\gamma\in E^{\Delta}$, as wanted. We have proven that $d^{*}$ is a fixpoint of $\Delta_{LP}$, and it is the greatest by Proposition 12.\par
 
% END A1.I2.ix2.p1
   
$\hfill\blacktriangleleft$\par
 
% END A1.SS0.SSS0.P0.SPx4.p1
 
% BEGIN A1.SS0.SSS0.P0.SPx4.p2
 
% BEGIN Thmtheorem16
 
\par\medskip\noindent\textbf{Proposition 16. [Restated, see original statement.]}\par
 
% BEGIN Thmtheorem16.p1
 
$\mathcal{D}_{LP}$ is a functor on $\mathbf{PMet_{1}}$.\par
 
% END Thmtheorem16.p1
 
% END Thmtheorem16
 
% END A1.SS0.SSS0.P0.SPx4.p2
 
% END A1.SS0.SSS0.P0.SPx4
